<%= heading2 %>{UML diagrams using PlantUML}

\href{https://plantuml.com}{PlantUML} lays out UML diagrams automatically; a relation is a single arrow such as \texttt{Car *- Wheel}.
It needs \texttt{lualatex}, the \texttt{--shell-escape} compiler flag, and the external PlantUML tool (which needs Java); the Docker image by DANTE e.V.\ contains all of this (see the \texttt{README}).
When these prerequisites are missing, the following demos are skipped automatically.

\ifluatex
  \ifnum\directlua{tex.sprint(ltg_uml_available and 1 or 0)}=1
    <%- bexample %>
    \begin{figure}
      \centering
      \begin{plantuml}
@startuml
Alice -> Bob: test
@enduml
\end{plantuml}
      \caption{Sequence diagram drawn with PlantUML}
      \label{fig:plantuml-sequence}
    \end{figure}
    <%- eexample %>

    <%- bexample %>
    \begin{figure}
      \centering
      \begin{plantuml}
@startuml
hide empty members
class Car
Driver - Car : drives >
Car *- Wheel : have 4 >
Car -- Person : < owns
@enduml
\end{plantuml}
      \caption{Class diagram drawn with PlantUML}
      \label{fig:plantuml-class-relations}
    \end{figure}
    <%- eexample %>

    <%- bexample %>
    \begin{figure}
      \centering
      \begin{plantuml}
@startuml
class BankAccount {
  +owner : String
  #balance : double
  -pin : int
  ~auditTrail : List
  +deposit(amount : double)
  #applyInterest()
  -validatePin(pin : int)
}
@enduml
\end{plantuml}
      \caption{Class with visibilities drawn with PlantUML}
      \label{fig:plantuml-class-visibility}
    \end{figure}
    <%- eexample %>

    <%- bexample %>
    \begin{figure}
      \centering
      \begin{plantuml}
@startuml
() "Interface 2" as I2
() "Interface 3" as I3
[component 1] as c1
[component 2] as c2
[component 3] as c3
c1 -- I2
c1 -- I3
I2 )-- c2
I3 )-- c3
@enduml
\end{plantuml}
      \caption{Component diagram drawn with PlantUML}
      \label{fig:plantuml-component}
    \end{figure}
    <%- eexample %>
  \else
    \emph{PlantUML or \texttt{--shell-escape} is not available, so these diagrams are not rendered. See the \texttt{README} on how to enable it.}
  \fi
\fi
